% ============================================================
% MỤC 3: PHÂN TÍCH VÀ THIẾT KẾ HỆ THỐNG
%   yêu cầu - use case - kiến trúc - giao thức - luồng dữ liệu - state machine - CSDL - bảo mật - giao diện - tính năng đặc biệt
% Phần nào chưa hiện thực thì ghi rõ "thiết kế" hoặc "đang phát triển".
% ============================================================
\section{Phân Tích Và Thiết Kế}

% ---------- Yêu cầu ----------
\begin{frame}{Yêu Cầu Chức Năng Và Phi Chức Năng}
  \begin{columns}[T]
    \begin{column}{0.52\textwidth}
      \begin{block}{Yêu cầu chức năng}
        \begin{enumerate}\setlength{\itemsep}{0pt}
          \item Đăng ký, đăng nhập, xem ai đang Online.
          \item Chat ở sảnh chờ.
          \item Tạo phòng, vào phòng bằng ID, tìm trận ngẫu nhiên.
          \item Xếp tàu, sẵn sàng (Ready), bắn theo lượt.
          \item Báo trúng/hụt, chìm tàu, kết thúc trận.
          \item Lưu lịch sử trận, xem bảng xếp hạng.
        \end{enumerate}
      \end{block}
    \end{column}
    \begin{column}{0.44\textwidth}
      \begin{alertblock}{Yêu cầu phi chức năng}
        \begin{itemize}\setlength{\itemsep}{0pt}
          \item \textbf{Đúng luật}: Server quyết định, chống gian lận.
          \item \textbf{Bảo mật}: mã hóa gói tin, băm mật khẩu, chống dò mật khẩu.
          \item \textbf{Đồng thời}: nhiều người chơi, không nghẽn.
          \item \textbf{Bền}: Client rớt mạng không làm Server sập.
          \item \textbf{Dễ chạy}: có chế độ SQLite offline để demo.
        \end{itemize}
      \end{alertblock}
    \end{column}
  \end{columns}
\end{frame}

% ---------- Use case ----------
\begin{frame}{Biểu Đồ Chức Năng Của Người Chơi}
  \centering
  \begin{tikzpicture}[
    grp/.style={draw=inkblue, fill=white, rounded corners=4pt, thick, align=left,
                font=\scriptsize, inner sep=6pt, text width=3.3cm},
    gt/.style={font=\scriptsize\bfseries, color=inkblue},
    arr/.style={-{Stealth[length=2mm]}, thick, color=brandblue}]
    \node[draw=inkblue, fill=gold!50, rounded corners=3pt, thick, minimum height=1.2cm, minimum width=1.6cm, font=\small\bfseries, align=center] (p) at (0,0) {Người\\chơi};
    \node[grp] (g1) at (4.3,1.9)  {\textbf{Tài khoản}\\Đăng ký\\Đăng nhập\\Xem Online/Offline};
    \node[grp] (g2) at (4.3,-1.9) {\textbf{Sảnh chờ}\\Chat sảnh\\Xem danh sách phòng\\Xem bảng xếp hạng};
    \node[grp] (g3) at (8.4,1.9)  {\textbf{Phòng và ghép trận}\\Tạo phòng\\Vào phòng bằng ID\\Tìm trận ngẫu nhiên\\Sẵn sàng (Ready)};
    \node[grp] (g4) at (8.4,-1.9) {\textbf{Trong trận}\\Xếp tàu\\Bắn theo lượt\\Dùng Radar\\Chat trong trận};
    \draw[arr] (p) -- (g1); \draw[arr] (p) -- (g2);
    \draw[arr] (g1) -- (g3); \draw[arr] (g2) -- (g4);
    \draw[arr] (g3) -- (g4);
  \end{tikzpicture}
  \par\vspace{0.1cm}\scriptsize Mũi tên giữa các nhóm thể hiện thứ tự dùng. Server là bên xử lý mọi yêu cầu.
\end{frame}

% ---------- Phân quyền ----------
\begin{frame}{Các Vai Trò Và Phân Quyền}
  \footnotesize
  \begin{tabularx}{\textwidth}{@{}Y c c c@{}}
    \toprule
    \textbf{Chức năng} & \textbf{Người chơi} & \textbf{Chủ phòng} & \textbf{Người vận hành} \\
    \midrule
    Đăng ký, đăng nhập, xem Online & $\checkmark$ & $\checkmark$ & \\
    Chat sảnh, xem bảng xếp hạng & $\checkmark$ & $\checkmark$ & \\
    Tạo phòng, vào phòng, tìm trận & $\checkmark$ & $\checkmark$ & \\
    Chơi trận (xếp tàu, bắn, Radar) & $\checkmark$ & $\checkmark$ & \\
    \textbf{Kick thành viên khỏi phòng} & & $\checkmark$ & \\
    Chạy Server, chọn cơ sở dữ liệu, xem log & & & $\checkmark$ \\
    Quản lý dữ liệu trên Supabase & & & $\checkmark$ \\
    \bottomrule
  \end{tabularx}
  \vspace{0.2cm}
  \begin{alertblock}{Ghi chú}
    \scriptsize \textbf{Chủ phòng} là người tạo phòng (hoặc người đứng đầu hàng đợi khi ghép trận), có thêm quyền kick; Server kiểm tra quyền này, Client không tự quyết. \textbf{Người vận hành} là nhóm phát triển, thao tác trực tiếp trên Server và Supabase. Hệ thống \textbf{không có tài khoản admin} trong phạm vi đề tài.
  \end{alertblock}
\end{frame}

% ---------- DFD mức 0 ----------
\begin{frame}{Lược Đồ Dòng Dữ Liệu (DFD) Mức 0}
  \centering
  \begin{tikzpicture}[
    ent/.style={draw=inkblue, fill=gold!40, thick, minimum width=2.6cm, minimum height=1cm, align=center, font=\scriptsize},
    proc/.style={draw=inkblue, fill=paleblue, thick, circle, minimum size=2.9cm, align=center, font=\small\bfseries},
    store/.style={draw=inkblue, fill=white, thick, minimum width=3cm, minimum height=0.9cm, align=center, font=\scriptsize},
    flow/.style={-{Stealth[length=2mm]}, thick, color=brandblue},
    lab/.style={font=\tiny, fill=bgblue, inner sep=1pt, align=center}]
    \node[proc] (p) at (0,0) {Hệ thống\\Battleship};
    \node[ent] (u) at (-5.6,0.6) {Người chơi};
    \node[ent] (o) at (-5.6,-1.8) {Người vận hành};
    \node[ent] (w) at (5.6,1.6) {Discord /\\Telegram};
    \node[store] (d) at (5.6,-0.6) {D: Cơ sở dữ liệu\\(Users, Matches, MatchDetails)};
    \draw[flow] ([yshift=4pt]u.east) -- node[lab, above]{yêu cầu: đăng nhập, phòng,\\tìm trận, bắn, chat} ([yshift=10pt]p.west);
    \draw[flow] ([yshift=-4pt]p.west) -- node[lab, below]{kết quả, trạng thái trận,\\danh sách Online} ([yshift=-4pt]u.east);
    \draw[flow] (o) -- node[lab, below]{khởi động, chọn CSDL} (p);
    \draw[flow] (p) -- node[lab, above]{ghi trận, đọc\\tài khoản} (d);
    \draw[flow] (p) -- node[lab, above, sloped]{kết quả trận (HTTP POST)} (w);
  \end{tikzpicture}
  \par\vspace{0.1cm}\scriptsize Người chơi còn đọc bảng xếp hạng trực tiếp qua REST (chỉ đọc). Webhook đang phát triển.
\end{frame}

% ---------- DFD mức 1 ----------
\begin{frame}{Lược Đồ Dòng Dữ Liệu (DFD) Mức 1}
  \centering
  \begin{tikzpicture}[
    proc/.style={draw=inkblue, fill=paleblue, thick, rounded corners=6pt, minimum width=2.9cm, minimum height=1.05cm, align=center, font=\scriptsize},
    ent/.style={draw=inkblue, fill=gold!40, thick, minimum width=1.9cm, minimum height=0.9cm, align=center, font=\scriptsize},
    store/.style={draw=inkblue, fill=white, thick, minimum width=2.8cm, minimum height=0.75cm, align=center, font=\scriptsize},
    flow/.style={-{Stealth[length=1.8mm]}, thick, color=brandblue},
    lab/.style={font=\tiny, fill=bgblue, inner sep=1pt}]
    \node[ent] (u) at (-5.8,0) {Người\\chơi};
    \node[proc] (p1) at (-2.2,2.1)  {\textbf{1} Xác thực};
    \node[proc] (p2) at (-2.2,0.35) {\textbf{2} Phòng và\\ghép trận};
    \node[proc] (p3) at (-2.2,-1.5) {\textbf{3} Điều khiển trận\\(đang phát triển)};
    \node[proc] (p4) at (-2.2,-3.2) {\textbf{4} Chat sảnh};
    \node[proc] (p5) at (2.1,-3.2)  {\textbf{5} Thống kê,\\bảng xếp hạng};
    \node[store] (d1) at (5.4,2.1)  {D1 Users};
    \node[store] (d2) at (5.4,-1.5) {D2 Matches,\\MatchDetails};
    \node[store] (d3) at (2.1,0.35) {D3 Phòng\\(trong bộ nhớ)};
    \draw[flow] (u) -- (p1); \draw[flow] (u) -- (p2); \draw[flow] (u) -- (p3); \draw[flow] (u) -- (p4);
    \draw[flow] (p1) -- node[lab, above]{tên, hash} (d1);
    \draw[flow] (p2) -- node[lab, above]{tạo, vào, rời} (d3);
    \draw[flow] (p3) -- node[lab, above]{ghi kết quả} (d2);
    \draw[flow] (p3) -- (d3);
    \draw[flow] (d1) to[bend left=15] node[lab, right]{Wins, Score} (p5);
    \draw[flow] (d2) -- (p5);
  \end{tikzpicture}
  \par\vspace{0.1cm}\scriptsize \textbf{1, 2, 4} đã chạy. \textbf{3} (luật trận) và ghi \textbf{D2} đang phát triển. \textbf{5} đọc qua view \texttt{leaderboard}.
\end{frame}

% ---------- Kiến trúc ----------
\begin{frame}{Kiến Trúc Tổng Thể}
  \centering
  \begin{tikzpicture}[
    box/.style={draw=inkblue, fill=white, rounded corners=3pt, thick, align=center,
                minimum height=1.05cm, font=\scriptsize},
    arr/.style={{Stealth[length=2mm]}-{Stealth[length=2mm]}, thick, color=brandblue},
    one/.style={-{Stealth[length=2mm]}, thick, color=brandblue, dashed}]
    \node[box, fill=paleblue, minimum width=2.6cm] (c) at (0,0) {\textbf{Client}\\WinForms\\\texttt{NetworkManager}};
    \node[box, fill=paleblue, minimum width=3.4cm] (s) at (5.2,0) {\textbf{Server}\\\texttt{TcpServer} $\rightarrow$ \texttt{PacketRouter}\\$\rightarrow$ Handler};
    \node[box, minimum width=2.6cm] (d) at (10.4,0.9) {\textbf{PostgreSQL}\\Supabase (cloud)};
    \node[box, minimum width=2.6cm] (l) at (10.4,-0.9) {\textbf{SQLite}\\(chạy offline)};
    \node[box, minimum width=2.6cm] (r) at (5.2,-2.5) {\textbf{REST (HTTPS)}\\view \texttt{leaderboard}};
    \node[box, minimum width=2.6cm] (w) at (5.2,2.4) {\textbf{Webhook}\\Discord/Telegram};
    \draw[arr] (c) -- node[above,font=\tiny, align=center]{TCP\\+ AES-128} (s);
    \draw[arr] (s) -- node[above,font=\tiny]{EF Core} (d);
    \draw[arr] (s) -- (l);
    \draw[one] (c) |- (r); \draw[one] (r) -| (d);
    \draw[one] (s) -- node[right,font=\tiny]{HTTP POST} (w);
  \end{tikzpicture}
  \par\vspace{0.1cm}
  \scriptsize Nét liền: đường chính (TCP, cơ sở dữ liệu). Nét đứt: chỉ đọc bảng xếp hạng qua REST; webhook (đang phát triển).
\end{frame}

% ---------- Gói tin ----------
\begin{frame}{Thiết Kế Giao Thức Gói Tin}
  \footnotesize
  \begin{columns}[T]
    \begin{column}{0.46\textwidth}
      \begin{block}{Cấu trúc một gói}
        \texttt{[4 byte độ dài][AES(JSON\{Type, Payload\})]}
        \begin{itemize}\setlength{\itemsep}{0pt}
          \item \texttt{Type} là kiểu liệt kê (\texttt{PacketType}).
          \item \texttt{Payload} là JSON của DTO tương ứng.
          \item Gói đầu tiên là Handshake (phiên bản giao thức).
        \end{itemize}
      \end{block}
    \end{column}
    \begin{column}{0.5\textwidth}
      \begin{tabularx}{\textwidth}{@{}lY@{}}
        \toprule \textbf{Nhóm} & \textbf{Kiểu gói chính} \\ \midrule
        Tài khoản & \texttt{LoginRequest/Response}, \texttt{RegisterRequest/Response} \\
        Phòng & \texttt{CreateRoomRequest}, \texttt{JoinRoomRequest/Response}, \texttt{RoomUpdate}, \texttt{LeaveRoom} \\
        Ghép trận & \texttt{MatchmakingJoin/CancelRequest}, \texttt{MatchFound} \\
        Trận đấu & \texttt{PlaceShipsRequest}, \texttt{PlayerReady}, \texttt{GameStart}, \texttt{FireShotRequest}, \texttt{ShotResult}, \texttt{TurnChange}, \texttt{GameOver} \\
        Khác & \texttt{ChatMessage}, \texttt{OnlineListUpdate}, \texttt{Ping/Pong}, \texttt{Error} \\
        \bottomrule
      \end{tabularx}
    \end{column}
  \end{columns}
\end{frame}

% ---------- Luồng đăng nhập ----------
\begin{frame}{Luồng Dữ Liệu: Đăng Nhập}
  \centering
  \begin{tikzpicture}[
    life/.style={dashed, color=skyblue, thick},
    msg/.style={-{Stealth[length=2mm]}, thick, color=brandblue},
    lab/.style={font=\scriptsize, fill=bgblue, inner sep=1pt}]
    \node[draw=inkblue, fill=paleblue, rounded corners=2pt, thick, minimum width=2cm, font=\small\bfseries] (c) at (0,0) {Client};
    \node[draw=inkblue, fill=paleblue, rounded corners=2pt, thick, minimum width=2cm, font=\small\bfseries] (s) at (5,0) {Server};
    \node[draw=inkblue, fill=paleblue, rounded corners=2pt, thick, minimum width=2cm, font=\small\bfseries] (d) at (10,0) {Cơ sở dữ liệu};
    \draw[life] (c.south) -- (0,-5.2); \draw[life] (s.south) -- (5,-5.2); \draw[life] (d.south) -- (10,-5.2);
    \draw[msg] (0,-1.0) -- node[lab, above]{1. Handshake (phiên bản)} (5,-1.0);
    \draw[msg] (5,-1.6) -- node[lab, above]{2. Chấp nhận} (0,-1.6);
    \draw[msg] (0,-2.3) -- node[lab, above]{3. \texttt{LoginRequest} (AES)} (5,-2.3);
    \draw[msg] (5,-3.0) -- node[lab, above]{4. Lấy hash theo tên} (10,-3.0);
    \draw[msg] (10,-3.6) -- node[lab, above]{5. Hash BCrypt} (5,-3.6);
    \node[font=\scriptsize, align=center, fill=gold!40, rounded corners=2pt] at (5,-4.2) {6. So khớp BCrypt, kiểm tra khóa và đăng nhập trùng};
    \draw[msg] (5,-4.9) -- node[lab, above]{7. \texttt{LoginResponse} (AES)} (0,-4.9);
  \end{tikzpicture}
\end{frame}

% ---------- Luồng một lượt bắn ----------
\begin{frame}{Luồng Dữ Liệu: Một Lượt Bắn (Thiết Kế)}
  \begin{columns}[T]
    \begin{column}{0.55\textwidth}
      \begin{block}{Server-Authoritative}
        \begin{enumerate}\setlength{\itemsep}{0pt}
          \item Client gửi \texttt{FireShotRequest} (tọa độ).
          \item Server kiểm tra: \textbf{đúng lượt?} ô \textbf{chưa bắn?} trận \textbf{đang chơi?} Sai thì bỏ qua.
          \item Server tra bàn của đối thủ, quyết định \textbf{Hit/Miss/Chìm}.
          \item Server gửi \texttt{ShotResult} cho \textbf{cả hai}, rồi \texttt{TurnChange}.
          \item Hết tàu: gửi \texttt{GameOver}, ghi kết quả vào cơ sở dữ liệu.
        \end{enumerate}
      \end{block}
      \begin{alertblock}{Đang phát triển}
        Logic trận trên Server (bước 2--5) chưa xong; xem lịch ở mục Giới thiệu.
      \end{alertblock}
    \end{column}
    \begin{column}{0.4\textwidth}
      \centering
      \begin{tikzpicture}[
        flow/.style={draw=inkblue, fill=white, rounded corners=3pt, thick, align=center,
                     minimum height=0.8cm, minimum width=4.2cm, font=\scriptsize},
        arr/.style={-{Stealth[length=2mm]}, thick, color=brandblue}]
        \node[flow, fill=paleblue] (a) at (0,0) {Client A: \texttt{FireShotRequest}};
        \node[flow] (b) at (0,-1.3) {Server kiểm tra lượt, ô, trạng thái};
        \node[flow] (c) at (0,-2.6) {Server tra bàn B, tính kết quả};
        \node[flow, fill=paleblue] (d) at (0,-3.9) {\texttt{ShotResult} gửi A và B};
        \draw[arr] (a) -- (b); \draw[arr] (b) -- (c); \draw[arr] (c) -- (d);
      \end{tikzpicture}
    \end{column}
  \end{columns}
\end{frame}

% ---------- State machine ----------
\begin{frame}{Máy Trạng Thái Của Phòng Và Trận}
  \begin{columns}[T]
    \begin{column}{0.46\textwidth}
      \begin{block}{Các trạng thái}
        \begin{itemize}\setlength{\itemsep}{0pt}
          \item \textbf{Waiting}: chờ đủ 2 người.
          \item \textbf{Placement}: hai bên xếp tàu.
          \item \textbf{Ready}: cả hai bấm sẵn sàng.
          \item \textbf{Playing}: bắn theo lượt.
          \item \textbf{GameOver}: có người thắng hoặc đối thủ mất kết nối.
        \end{itemize}
      \end{block}
      \begin{exampleblock}{Quy tắc}
        Chỉ nhận lệnh bắn khi ở \textbf{Playing} và đúng lượt. Bắt đầu trận chỉ khi \textbf{cả hai} Ready.
      \end{exampleblock}
    \end{column}
    \begin{column}{0.5\textwidth}
      \centering
      % [ẢNH] sơ đồ State Machine Đức đã vẽ ở README
      \slideimg[\textwidth]{12-state-machine.png}{Hình State Machine của Đức (lưu từ README trên GitHub)}
    \end{column}
  \end{columns}
\end{frame}

% ---------- ERD ----------
\begin{frame}{Thiết Kế Cơ Sở Dữ Liệu (ERD)}
  \centering
  \begin{tikzpicture}[
    ent/.style={draw=inkblue, fill=white, rounded corners=2pt, thick, align=left,
                font=\scriptsize, inner sep=5pt},
    arr/.style={thick, color=brandblue}]
    \node[ent] (u) at (0,0)    {\textbf{Users}\\[2pt]\underline{Id} (PK)\\Username (UNIQUE)\\Password (BCrypt)\\Wins, Losses, Score\\Condition\\UserId (FK, bạn bè)};
    \node[ent] (m) at (5.4,0)  {\textbf{Matches}\\[2pt]\underline{Id} (PK)\\Player1Id (FK)\\Player2Id (FK)\\WinnerId (FK)\\StartTime, EndTime};
    \node[ent] (d) at (10.8,0) {\textbf{MatchDetails}\\[2pt]\underline{Id} (PK)\\MatchId (FK)\\PlayerId (FK)\\X\_Coordinate, Y\_Coordinate\\IsHit, TurnTime};
    \draw[arr] (u) -- node[above,font=\scriptsize]{1 -- N} (m);
    \draw[arr] (m) -- node[above,font=\scriptsize]{1 -- N} (d);
  \end{tikzpicture}

  \vspace{0.3cm}\footnotesize
  Một người chơi có nhiều trận; một trận có nhiều lượt bắn (\textbf{MatchDetails}). View \texttt{leaderboard} chỉ lấy tên, thắng, thua, điểm từ \texttt{Users}.
\end{frame}

% ---------- Bảo mật ----------
\begin{frame}{Thiết Kế Bảo Mật Nhiều Lớp}
  \footnotesize
  \begin{tabularx}{\textwidth}{@{}l Y Y@{}}
    \toprule \textbf{Lớp} & \textbf{Biện pháp} & \textbf{Chống lại} \\ \midrule
    Đường truyền & AES-128-CBC, IV ngẫu nhiên mỗi gói & Nghe lén, đọc trộm gói tin \\
    Mật khẩu & BCrypt (có salt), kiểm tra dù tên không tồn tại & Lộ mật khẩu nếu rò rỉ DB, dò tên tài khoản \\
    Đăng nhập & Khóa 60 giây sau 5 lần sai (theo IP và theo tên) & Dò mật khẩu hàng loạt \\
    Đăng ký & Tối đa 5 lần mỗi 10 phút mỗi IP & Tạo tài khoản rác \\
    Phiên & Router chặn mọi gói (trừ đăng nhập/đăng ký) khi chưa đăng nhập; một tài khoản một phiên & Gửi lệnh khi chưa xác thực, đăng nhập trùng \\
    Dữ liệu & RLS; khóa công khai chỉ đọc view bảng xếp hạng & Đọc trực tiếp bảng \texttt{Users} từ ngoài \\
    Luật chơi & Server quyết định Hit/Miss & Client gian lận \\
    \bottomrule
  \end{tabularx}
\end{frame}

% ---------- Giao diện ----------
\begin{frame}{Thiết Kế Giao Diện}
  \centering
  \begin{columns}[T]
    \begin{column}{0.31\textwidth}
      \centering\scriptsize\textbf{Đăng nhập}\par
      \slideimg[\textwidth]{04-client-login.png}{Màn hình Đăng nhập}
    \end{column}
    \begin{column}{0.31\textwidth}
      \centering\scriptsize\textbf{Sảnh chờ (Lobby)}\par
      \slideimg[\textwidth]{05-client-lobby.png}{Màn hình Lobby sau khi đăng nhập}
    \end{column}
    \begin{column}{0.31\textwidth}
      \centering\scriptsize\textbf{Bảng xếp hạng}\par
      \slideimg[\textwidth]{06-client-leaderboard.png}{Form Bảng xếp hạng có dữ liệu}
    \end{column}
  \end{columns}
  \vspace{0.2cm}
  \begin{columns}[T]
    \begin{column}{0.31\textwidth}
      \centering\scriptsize\textbf{Đăng ký}\par
      \slideimg[\textwidth]{gk-02-register.png}{Màn hình Đăng ký}
    \end{column}
    \begin{column}{0.31\textwidth}
      \centering\scriptsize\textbf{Phòng chờ}\par
      \slideimg[\textwidth]{gk-03-roomwait.png}{Phòng chờ có 2 người chơi và nút Ready}
    \end{column}
    \begin{column}{0.31\textwidth}
      \centering\scriptsize\textbf{Bàn cờ}\par
      \slideimg[\textwidth]{gk-04-board.png}{Lưới 10x10 (chụp từ Client)}
    \end{column}
  \end{columns}
\end{frame}

% ---------- Tính năng đặc biệt ----------
\begin{frame}{Thiết Kế Tính Năng Đặc Biệt}
  \begin{columns}[T]
    \begin{column}{0.48\textwidth}
      \begin{block}{Radar (quét vùng 3x3)}
        \begin{itemize}\setlength{\itemsep}{0pt}
          \item Mỗi người dùng \textbf{một lần mỗi trận}, chọn một ô.
          \item Server đếm số ô tàu trong vùng $3\times3$ quanh ô đó và chỉ trả về \textbf{con số}, không nói ô nào.
          \item Dùng Radar thì \textbf{mất một lượt}: phải cân nhắc giữa quét và bắn.
          \item Server giữ cờ ``đã dùng'' và áp luật đúng lượt: Client không dùng được hai lần.
        \end{itemize}
      \end{block}
    \end{column}
    \begin{column}{0.48\textwidth}
      \begin{block}{Gửi kết quả trận ra Discord/Telegram}
        \begin{itemize}\setlength{\itemsep}{0pt}
          \item Khi trận kết thúc, Server gọi \textbf{HTTP POST} tới webhook.
          \item Tin nhắn xuất hiện ngay trên điện thoại: ``A thắng B sau N lượt''.
          \item Nối mạng game (TCP) với Internet (HTTP) và dịch vụ bên ngoài.
          \item Chạy ở Task riêng với thời hạn 5 giây; lỗi chỉ ghi log, không làm treo Server. URL webhook để trong biến môi trường.
        \end{itemize}
      \end{block}
    \end{column}
  \end{columns}
  \vspace{0.1cm}
  \begin{alertblock}{Trạng thái}
    \footnotesize Đang ở bước thiết kế; hiện thực sau khi Server có logic trận (dự kiến 22/11 -- 25/11).
  \end{alertblock}
\end{frame}
